\documentclass[12pt,fleqn,twoside,a4paper]{report}
\usepackage{etex}
% Select encoding of your inputs.
\usepackage[utf8]{inputenc}

% Make latex understand and use the typographic
% rules of the language used in the document.
\usepackage[english, danish]{babel}

% Use the vector font Latin Modern which is going
% to be the default font in latex in the future.
\usepackage{lmodern}
% Use vector arrows in mathmode
\usepackage{esvect}

% Choose the font encoding
\usepackage[T1]{fontenc}

% load a colour package
\usepackage[table]{xcolor}
\definecolor{aaublue}{RGB}{33,26,82}% dark blue

% The standard graphics inclusion package
\definecolor{white}{RGB}{255,255,255} % define color white
\usepackage{graphicx}

% Set up how figure and table captions are displayed
\usepackage{caption}
\captionsetup{%
  font=footnotesize,% set font size to footnotesize
  labelfont=bf % bold label (e.g., Figure 3.2) font
}

% Enable row combination in tables
\usepackage{multirow}

% Make space between table lines and text
\renewcommand{\arraystretch}{1.5}

% Make the standard latex tables look so much better
\usepackage{array,booktabs}

% Enable the use of frames around, e.g., theorems
% The framed package is used in the example environment
\usepackage{framed}
\usepackage{colortbl}
\usepackage{longtable}

\usepackage{textcomp}

%%%%%%%%%%%%%%%%%%%%%%%%%%%%%%%%%%%%%%%%%%%%%%%%
% Mathematics
%%%%%%%%%%%%%%%%%%%%%%%%%%%%%%%%%%%%%%%%%%%%%%%%
% Defines new environments such as equation,
% align and split 
\usepackage{amsmath}
\usepackage{relsize}
% Adds new math symbols
\usepackage{amssymb}
% Use theorems in your document
% The ntheorem package is also used for the example environment
% When using thmmarks, amsmath must be an option as well. Otherwise \eqref doesn't work anymore.
\usepackage[framed,amsmath,thmmarks]{ntheorem}

%%%%%%%%%%%%%%%%%%%%%%%%%%%%%%%%%%%%%%%%%%%%%%%%
% Page Layout
%%%%%%%%%%%%%%%%%%%%%%%%%%%%%%%%%%%%%%%%%%%%%%%%
% Change margins, papersize, etc of the document
\usepackage[
  left=25mm,% left margin on an odd page %tidligere 25mm for baade right og left
  right=25mm,% right margin on an odd page
  top=35mm,
  ]{geometry}
  
% Modify how \chapter, \section, etc. look
% The titlesec package is very configureable
\usepackage{titlesec}
\makeatletter
\def\ttl@mkchap@i#1#2#3#4#5#6#7{%
    \ttl@assign\@tempskipa#3\relax\beforetitleunit
    \vspace{\@tempskipa}%<<<<<< REMOVE THE * AFTER \vspace
    \global\@afterindenttrue
    \ifcase#5 \global\@afterindentfalse\fi
    \ttl@assign\@tempskipb#4\relax\aftertitleunit
    \ttl@topmode{\@tempskipb}{%
        \ttl@select{#6}{#1}{#2}{#7}}%
    \ttl@finmarks  % Outside the box!
    \@ifundefined{ttlp@#6}{}{\ttlp@write{#6}}}
\makeatother

\titlespacing{\chapter}{0pt}{0pt}{10pt}
\titlespacing{\section}{0pt}{0pt}{-5pt}
\titlespacing{\subsection}{0pt}{8pt}{-5pt}
\titlespacing{\subsubsection}{0pt}{6pt}{-10pt}

\titleformat*{\section}{\normalfont\Large\bfseries}%\color{aaublue}}
\titleformat*{\subsection}{\normalfont\large\bfseries}%\color{aaublue}}
\titleformat*{\subsubsection}{\normalfont\normalsize\bfseries}%\color{aaublue}}
%\titleformat*{\paragraph}{\normalfont\normalsize\bfseries\color{aaublue}}
%\titleformat*{\subparagraph}{\normalfont\normalsize\bfseries\color{aaublue}}

%Formatting chapter headers. Predefined styles that can be used as option to the package: Sonny, Lenny, Glenn, Conny, Rejne, Bjarne, PetersLenny and Bjornstrup

%\usepackage[Sonny]{fncychap}

\usepackage{titlesec, blindtext, color}
%\definecolor{gray75}{gray}{0.75}
\newcommand{\hsp}{\hspace{20pt}}
\titleformat{\chapter}[hang]{\Huge\bfseries}{\thechapter\hsp{|}\hsp}{0pt}{\Huge\bfseries}


% Change the headers and footers
\usepackage{fancyhdr}
\setlength{\headheight}{15pt}
\pagestyle{fancy}
\fancyhf{} %delete everything
\renewcommand{\headrulewidth}{0pt} %remove the horizontal line in the header
\fancyhead[RO,LE]{\small\nouppercase\leftmark} %even page - chapter title

\fancyhead[LO]{}
\fancyhead[RE]{} 
\fancyhead[CE]{}
\fancyhead[CO]{}

\fancyfoot[RE,LO]{\thepage}
\fancyfoot[LE,RO]{B205} %page number on all pages
\fancyfoot[CE,CO]{}



% change first page of all chapters header and footer to fancy style
\makeatletter
\let\ps@plain\ps@fancy
\makeatother

% Do not stretch the content of a page. Instead,
% insert white space at the bottom of the page
\raggedbottom

% Enable arithmetics with length. Useful when typesetting the layout.
\usepackage{calc}

%%%%%%%%%%%%%%%%%%%%%%%%%%%%%%%%%%%%%%%%%%%%%%%%
% Bibliography
%%%%%%%%%%%%%%%%%%%%%%%%%%%%%%%%%%%%%%%%%%%%%%%%
% Add the \citep{key} command which display a
% reference as [author, year]
\usepackage[square]{natbib}
\usepackage{natbib}
% Appearance of the bibliography
\bibliographystyle{00setup/apalike}
%\bibliographystyle{IEEEtran}

%%%%%%%%%%%%%%%%%%%%%%%%%%%%%%%%%%%%%%%%%%%%%%%%
% Misc
%%%%%%%%%%%%%%%%%%%%%%%%%%%%%%%%%%%%%%%%%%%%%%%%

% Add bibliography and index to the table of
% contents
\usepackage[nottoc, notlof, notlot]{tocbibind}
\usepackage{tocloft}
% Add the command \pageref{LastPage} which refers to the
% page number of the last page
\setlength{\cftbeforetoctitleskip}{0 cm}

\usepackage[
  %disable, %turn off todonotes
  colorinlistoftodos, %enable a coloured square in the list of todos
  %inline,
  textwidth=20mm, %set the width of the todonotes
  textsize=scriptsize, %size of the text in the todonotes
  ]{todonotes}
\setlength{\marginparwidth}{2cm}
\newcommand{\todocheck}[1]{\todo[color=green!40, author=Check]{#1}}
\newcommand{\todojens}[1]{\todo[color=red!40, author=Jens, inline]{#1}}
%Command for inserting green todo notes as Ole's corrections.
%\renewcommand{\todo}[1]{\todo[noline]{#1}}  
    
\usepackage{wrapfig}
\usepackage{amssymb}
\usepackage{pifont}
% Enables figures with text wrapped tightly around it 

%%% Table of contents %%%%%  
\setcounter{tocdepth}{1}

\renewcommand{\cftpartpresnum}{Part~}
\let\cftoldpartfont\cftpartfont
\renewcommand{\cftpartfont}{\cftoldpartfont\cftpartpresnum}

%%%%%%%%%%%%%%%%%%%%%%%%%%%%%%%%%%%%%%%%%%%%%%%%
% Hyperlinks
%%%%%%%%%%%%%%%%%%%%%%%%%%%%%%%%%%%%%%%%%%%%%%%%

% Enable hyperlinks and insert info into the pdf
% file. Hypperref should be loaded as one of the 
% last packages
\usepackage{nameref}
\usepackage{hyperref}
\hypersetup{%
	%pdfpagelabels=true,%
	plainpages=false,%
	pdfauthor={Author(s)},%
	pdftitle={Title},%
	pdfsubject={Subject},%
	bookmarksnumbered=true,%
	colorlinks,%
	citecolor=black,%
	filecolor=black,%
	linkcolor=black,% you should probably change this to black before printing
	urlcolor=black,%
	pdfstartview=FitH%
}

% remove all indentations
\setlength\parindent{0pt}
\parskip 5mm
\usepackage{verbatim}


\definecolor{Gra}{RGB}{230,230,230}

%creates a nice-looking C#-text
\newcommand{\CC}{C\nolinebreak\hspace{-.05em}\raisebox{.3ex}{\scriptsize\text \#} }

%enables multi column lists
\usepackage{multicol}

%enables code-examples
\usepackage{listings}

\definecolor{coolblue}{RGB}{32,95,128}
\definecolor{mygreen}{rgb}{0,0.6,0}
\definecolor{mygray}{rgb}{0.5,0.5,0.5}
\definecolor{mymauve}{rgb}{0.58,0,0.82}
\usepackage{textcomp}
\definecolor{listinggray}{gray}{0.9}
\definecolor{lbcolor}{rgb}{0.9,0.9,0.9}


%%%%%%%%%%%%%%%%%%%%%%%%%%%%%%%%%%%%%%%%%%%%%%%%
% Listing
%%%%%%%%%%%%%%%%%%%%%%%%%%%%%%%%%%%%%%%%%%%%%%%%

\usepackage{xcolor}
\colorlet{keyword}{blue!90!black!100}
\colorlet{keyword2}{red!40!blue!85}
\colorlet{comment}{green!60!black!100}
\definecolor{backcolour}{rgb}{0,0,0}
\usepackage{beramono}% monospaced font with bold variant
\definecolor{backcolourlst}{rgb}{0.97,0.97,0.96}
%\definecolor{backcolourlst}{rgb}{1,0.75,0.80}

\lstset{
%  backgroundcolor=\color{white},   % choose the background color; you must add \usepackage{color} or \usepackage{xcolor}
%  basicstyle=\footnotesize,        % the size of the fonts that are used for the code
%  breakatwhitespace=false,         % sets if automatic breaks should only happen at whitespace
%  breaklines=true,                 % sets automatic line breaking
%  captionpos=t,                    % sets the caption-position to bottom
%  commentstyle=\color{mygreen},    % comment style
%  deletekeywords={...},            % if you want to delete keywords from the given language
%  escapeinside={\%*}{*)},          % if you want to add LaTeX within your code
%  extendedchars=true,              % lets you use non-ASCII characters; for 8-bits encodings only, does not work with UTF-8
%  frame=single,                    % adds a frame around the code
%  keepspaces=true,                 % keeps spaces in text, useful for keeping indentation of code (possibly needs columns=flexible)
%  keywordstyle=\color{blue},       % keyword style
%  language=C++,                 % the language of the code
%  morekeywords={*,...},            % if you want to add more keywords to the set
%  numbers=left,                    % where to put the line-numbers; possible values are (none, left, right)
%  numbersep=5pt,                   % how far the line-numbers are from the code
%  numberstyle=\tiny\color{mygray}, % the style that is used for the line-numbers
%  rulecolor=\color{black},         % if not set, the frame-color may be changed on line-breaks within not-black text (e.g. comments (green here))
%  showspaces=false,                % show spaces everywhere adding particular underscores; it overrides 'showstringspaces'
%  showstringspaces=false,          % underline spaces within strings only
%  showtabs=false,                  % show tabs within strings adding particular underscores
%  stepnumber=1,                    % the step between two line-numbers. If it's 1, each line will be numbered
%  stringstyle=\color{mymauve},     % string literal style
%  tabsize=2,                       % sets default tabsize to 2 spaces
%  title=\lstname                   % show the filename of files included with \lstinputlisting; also try caption instead of title
	backgroundcolor=\color{lbcolor},
	tabsize=4,
	rulecolor=\color{black},
	language=C,
  	basicstyle   = \linespread{1.2} \ttfamily \scriptsize,
  	keywordstyle = \color{keyword}\bfseries,	
  	commentstyle = \color{comment}\bfseries,
  	numbers=left,                   
	numberstyle=\scriptsize,      
	stepnumber=1,                   
	numbersep=8pt,                  
	backgroundcolor=\color{white},  
	showspaces=false,               
	showstringspaces=false,         
	showtabs=false,                 
	backgroundcolor=\color{backcolourlst},
	%frame=single,
	frame=single,                   
	captionpos=b,           
	breaklines=true,        
	breakatwhitespace=false,    
	escapeinside={\%*}{*)}
}
\lstset{
	emph={uint8, int8_t, uint16_t, int16, uint32_t, int32_t,}, emphstyle={\color{keyword}\bfseries},
    emph={[2]ISR, TCE1, CNT, TCC1, TCD1_OVF_vect},emphstyle={[2]\color{keyword2}\bfseries}%
}



\usepackage{float}
\usepackage{caption}
\usepackage{subcaption}
\usepackage{siunitx}
\sisetup{decimalsymbol=comma}
\sisetup{detect-weight}

\usepackage{enumitem}
%\usepackage[citestyle=authoryear,natbib=true]{biblatex}

% Figures - TIKZ
\usepackage{tikz}
\usepackage[americanresistors,americaninductors,americancurrents, americanvoltages]{circuitikz}


% Wall of text logo

\newcommand{\walloftextalert}[0]{\includegraphics[width=\textwidth]{walloftext.png}}


% Citation aliasses for nicer references in text
\defcitealias{ISO226}{\scshape ISO223, 2003}
\defcitealias{radio_standard}{\scshape IEC 61305-2, 1997}
\defcitealias{tysk}{\scshape EN 60315-4, 1998}
\defcitealias{REC}{\scshape Rec. ITU-R BS.704, 1990}
\defcitealias{std581}{\scshape IEC 581-6, 1979}
\defcitealias{std1305}{\scshape IEC 1305-3, 1995}
\defcitealias{std268-3}{\scshape IEC 60268-3, 2000}
\defcitealias{DS61938}{\scshape DS/EN 61938, 1997}
\defcitealias{std581-8}{\scshape IEC 581-8, 1987}


% URL break fix nede i bibfil

\def\UrlBreaks{\do\/\do-}


\usepackage{pdfpages}

\definecolor{blockcolor}{RGB}{112,168,218}

\tikzstyle{blockbig} = [draw, fill=blockcolor, rectangle, 
    minimum height=3em, minimum width=6em]
\tikzstyle{block} = [draw, fill=blockcolor, rectangle, 
    minimum height=2em, minimum width=2em]
    
\tikzstyle{blockbigclear} = [draw, fill=white, rectangle, 
    minimum height=3em, minimum width=6em]
\tikzstyle{blockclear} = [draw, fill=white, rectangle, 
    minimum height=2em, minimum width=2em]
    
\tikzstyle{sum} = [draw, fill=blockcolor, circle, node distance=1.5cm, minimum size=0.6cm]
\tikzstyle{sumclear} = [draw, fill=white, circle, node distance=1.5cm, minimum size=0.6cm]
\tikzstyle{input} = [coordinate]
\tikzstyle{output} = [coordinate]
\tikzstyle{pinstyle} = [pin edge={to-,thin,black}]

%% Centering columns of custom width
\newcolumntype{P}[1]{>{\centering\arraybackslash}p{#1}}
\newcolumntype{M}[1]{>{\centering\arraybackslash}m{#1}}

\usepackage{empheq}
\newcommand*\widefbox[1]{\fbox{\hspace{2em}#1\hspace{2em}}}

\newcommand{\todoanders}[1]{\todo[color=red!40, author=Anders]{#1}}
\graphicspath{{./rapport/billeder/}}
\newcommand{\todohusk}[1]{\todo[color=red, author=HUSK:]{#1}}


\makeatletter
\renewcommand{\todo}[2][]{\tikzexternaldisable\@todo[#1]{#2}\tikzexternalenable}
\renewcommand{\missingfigure}[2][]{\tikzexternaldisable\@missingfigure[#1]{#2}\tikzexternalenable}
\makeatother


%%PGF TABLES%%
\usepackage{booktabs}
\newcommand{\ra}[1]{\renewcommand{\arraystretch}{#1}}
\usepackage{array}
\usepackage{ragged2e}

\newcommand*{\textlabel}[2]{				% Gør \textlabel mulig
  \edef\@currentlabel{#1}					% Set target label
  \phantomsection							% Correct hyper reference link
  #1\label{#2}}								% Print and store label

% ¤¤ Opsætning af enheder ¤¤ % 
\newcommand{\enhed}[1]{\hfill[\mathrm{#1}]}   % Skriv \enhed{"din enhed"} og dem bliver rykket til højre




